% Generated mechanically from frozen input p12/ko/desirables.tex.
% Do not edit this generated file; edit the bound locale source instead.

% OK, start here.
%

\setcounter{chapter}{112}
\chapter{희망 사항}


\phantomsection
\label{desirables-section-phantom}



\section{서론}
\label{desirables-section-introduction}

\noindent
이 장은 기본적으로 스택 프로젝트에 넣고 싶은 것들을 나열한 목록일 뿐이다.
스택 프로젝트에 자료를 계속 추가하고 있으므로, 이 목록은 언제나
스택 프로젝트의 현재 상태보다 어느 정도 뒤처져 있다. 사실 추가해야 할
것들을 목록으로 만들려 한 것 자체가 실수였을지도 모른다. 이 목록을
최신 상태로 유지하는 것은 불가능해 보이기 때문이다.

\medskip\noindent
마지막 갱신: 2017년 8월 31일 목요일.


\section{규약}
\label{desirables-section-conventions}

\noindent
이 문서에서 사용하는 규약을 짧게 나열한 장이 있어야 한다.
그러한 장은 이미 존재한다. 『규약』의 절
\ref{conventions-section-comments}을 보라. 하지만 그곳에는 훨씬 더 많은
내용을 추가할 수 있다. 특히 ``숨겨진'' 규약과 암묵적인 가정을 찾아
그곳에 기록하면 유용할 것이다.


\section{사이트와 토포스}
\label{desirables-section-sites}

\noindent
사이트와 층을 다루는 장이 있다. 『사이트와 층』의 절
\ref{sites-section-introduction}을 보라.
환 달린 사이트(및 토포스)와 그 위의 가군을 다루는 장도 있다.
『사이트 위의 가군』의 절 \ref{sites-modules-section-introduction}을 보라.
이러한 상황에서의 코호몰로지를 다루는 장도 있다. 『사이트 코호몰로지』의
절 \ref{sites-cohomology-section-introduction}을 보라.
하지만 특히 코호몰로지를 다루는 장에는 훨씬 더 많은 내용을 추가할 수 있다.


\section{스택}
\label{desirables-section-stacks}

\noindent
(추상적) 스택을 다루는 장이 있다. 『스택』의 절
\ref{stacks-section-introduction}을 보라.
다음과 같이 개선하면 좋을 것이다.
\begin{enumerate}
\item ``스택화''에 관한 논의를 개선한다.
\item 스택화의 예를 제시한다.
\item 전반적으로 더 많은 예를 제시한다.
\item 거브에 관한 논의를 개선한다.
\end{enumerate}
아직 추가되지 않은 결과의 예를 하나 들자. 아벨 군층 $\mathcal{F}$가
$\mathcal{C}$ 위에 주어졌을 때, $\mathcal{F}$를 밴드로 갖는 거브들의
동치류 집합은 $H^2(\mathcal{C}, \mathcal{F})$와 일대일 대응한다.


\section{단체적 방법}
\label{desirables-section-simplicial}

\noindent
단체적 방법을 다루는 장이 있다. 『단체적 방법』의 절
\ref{simplicial-section-introduction}을 보라.
이 장은 검토하고 개선해야 한다. 단체적 호모토피(조합적 호모토피라고도 한다)와
Kan 복합체 사이의 관계에 관한 논의를 개선해야 한다.
단체적 공간을 다루는 장도 있다. 『단체적 대수공간』의 절
\ref{spaces-simplicial-section-introduction}을 보라.
이 장에서는 단체적 위상 공간, 단체적 사이트, 단체적 토포스를 간략히 논한다.
``단체적 대수기하학''을 더 발전시켜 단체적 스킴(또는 단체적 대수공간이나
단체적 대수적 스택)을 논하고, 기하학적 문제와 그 코호몰로지 등을 다룰 수 있다.


\section{스킴의 코호몰로지}
\label{desirables-section-schemes-cohomology}

\noindent
준연접 층의 코호몰로지를 다루는 장은 이미 있다. 『스킴의 코호몰로지』의
절 \ref{coherent-section-introduction}을 보라.
스킴 위의 준연접 층의 유도 범주를 논하는 장도 있다. 『스킴의 유도 범주』의
절 \ref{perfect-section-introduction}을 보라.
뇌터 스킴의 쌍대성과 스킴의 사상에 대한 상대 쌍대성을 논하는 장도 있다.
『스킴의 쌍대성』의 절 \ref{duality-section-introduction}을 보라.
또한 스킴의 에탈 코호몰로지와 결정 코호몰로지를 다루는 장들도 있다.
그러나 이 장들에 있는 내용은 대부분 매우 기초적이며, 훨씬 더 많은 내용을
추가할 수 있고 또 추가해야 한다.


\section{Schlessinger 방식의 변형 이론}
\label{desirables-section-deformation-schlessinger}

\noindent
이 내용을 다루는 장이 있다. 『형식 변형 이론』의 절
\ref{formal-defos-section-introduction}을 보라.
일반 이론의 예들을 논하는 장도 있다. 『변형 문제』의 절
\ref{examples-defos-section-introduction}을 보라.
또한 『변형 이론』의 절 \ref{defos-section-introduction}에서는
환(및 가군)의 변형, 환 달린 공간(및 가군층)의 변형,
환 달린 토포스(및 가군층)의 변형을 논한다.
이 장에서는 단순 코탄젠트 복합체를 사용하여 장애, 일차 변형,
무한소 자기동형을 기술한다. 이 내용은 뒤쪽 장들에서 모듈라이 스택의
대수성을 밝히는 데 몇 차례 응용되었다.
완전한 코탄젠트 복합체를 논하는 장도 있다. 『코탄젠트 복합체』의 절
\ref{cotangent-section-introduction}을 보라.


\section{대수적 스택의 정의}
\label{desirables-section-definition-algebraic-stacks}

\noindent
대수적 스택이란 fppf 위상을 갖춘 스킴의 범주 위의 준군 스택으로서,
대각사상이 대수공간으로 표현가능하고 어떤 스킴에서 오는 전사 매끄러운
사상의 공역이 되는 것을 말한다. 『대수적 스택』의 절
\ref{algebraic-section-algebraic-stacks}을 보라.
``Deligne-Mumford 스택''이란 어떤 스킴과 그 스킴에서 해당 스택으로 가는
전사 에탈 사상이 존재하는 대수적 스택을 말한다. 이는 Deligne과 Mumford의
논문 \cite{DM}에서와 같다. 『대수적 스택』의 정의
\ref{algebraic-definition-deligne-mumford}를 보라.
``Artin 스택''이라는 용어는 Artin의 논문들에 나오는 종류의 스택을 가리키는
데 쓰기로 한다. \cite{ArtinI}, \cite{ArtinII}, \cite{ArtinVersal}을 보라.
한 가지 가능한 정의는, 대수적 스택 $\mathcal{X}$가 국소 뇌터 스킴
$S$ 위에 있고 $\mathcal{X} \to S$가 국소적으로 유한형인 것\footnote{즉,
이들은 정확히 『Artin의 공리』의 절 \ref{artin-section-axioms}에 나오는
Artin의 공리 [-1], [0], [1], [2], [3], [4], [5]를 만족하는 $S$ 위의
대수적 스택들이다.}을 Artin 스택이라 하는 것이다.


\section{스킴, 대수공간, 대수적 스택의 예}
\label{desirables-section-examples-stacks}

\noindent
현재 스택 프로젝트에는 모듈라이 스택과 그 성질을 논하는 장이 두 개 있다.
『모듈라이 스택』의 절 \ref{moduli-section-introduction}과
『곡선의 모듈라이』의 절 \ref{moduli-curves-section-introduction}을 보라.
시간이 지나면서 다음과 같은 예를 더 추가하려 한다.
\begin{enumerate}
\item $\mathcal{A}_g$, 즉 종수 $g$인 주극성 아벨 스킴,
\item $\mathcal{A}_1 = \mathcal{M}_{1, 1}$, 즉 $1$점이 표시된
매끄러운 사영 종수 $1$ 곡선,
\item $\mathcal{M}_{g, n}$, 즉 매끄러운 사영 종수 $g$ 곡선으로서
서로 다른 표지된 점 $n$개를 갖는 것,
\item $\overline{\mathcal{M}}_{g, n}$, 즉 점 $n$개가 표시된 안정적인
마디 사영 종수 $g$ 곡선,
\item $\SheafHom_S(\mathcal{X}, \mathcal{Y})$, 즉 사상들의 모듈라이
(스택 $\mathcal{X}$, $\mathcal{Y}$와 바탕 스킴 $S$에 적절한 조건을 둔다),
\item $\textit{Bun}_G(X) = \SheafHom_S(X, BG)$, 즉 기하 Langlands
프로그램의 $G$-다발 스택(스킴 $X$, 군 스킴 $G$, 바탕 스킴 $S$에
적절한 조건을 둔다),
\item $\Picardstack_{\mathcal{X}/S}$, 즉 바탕 스킴(또는 공간) 위의
대수적 스택에 결부된 Picard 스택.
\end{enumerate}
더 일반적으로 말해, 스택 프로젝트에는 기하학적으로 의미 있는 예가
다소 부족하다.


\section{대수적 스택의 성질}
\label{desirables-section-stacks-properties}

\noindent
기본 이론이 이제 대부분 갖추어졌으므로, 이는 아마 작업하기 더 쉬운
프로젝트 가운데 하나일 것이다. 물론 이러한 것들은 실제로 스택의 사상의
성질이다. 매끄러운 인자를 무시하고 특이점 등을 정의할 수 있다.
연결된 정규 스택이 기약임을 증명하는 등의 일을 할 수 있다.


\section{대수적 스택의 lisse-\'etale 사이트}
\label{desirables-section-lisse-etale}

\noindent
이는 『대수적 스택의 코호몰로지』의 절
\ref{stacks-cohomology-section-lisse-etale}에서 도입되었다.
대수적 스택의 $1$-사상에 대해 함자적이지 않음을 보여 주는 예는
『예』의 절 \ref{examples-section-lisse-etale-not-functorial}에서 논한다.
물론 이에 관해 훨씬 더 많은 이야기를 할 수 있지만, 가능한 한
``큰'' 에탈 사이트를 사용하여 명제를 증명하는 것이 매우 유용하다는 사실이
드러난다.



\section{늘 알고 싶었지만 감히 묻지 못했던 것들}
\label{desirables-section-stacks-fun-lemmas}

\noindent
되풀이해서 사용하는 유용한 보조정리 가운데 문헌에서 따로 언급되지 않거나
참고문헌을 찾기 쉽지 않은 것들이 아주 많을 것이다. 기법 모음이다.

\medskip\noindent
예: 스킴의 두 준군 $R\Rightarrow U$와 $R' \Rightarrow U'$가 주어졌을 때,
스킴의 준군만을 사용하여 $1$-사상 $[U/R] \to [U'/R']$을 갖는다는 것은
정확히 무슨 뜻인가?



\section{스택 위의 준연접 층}
\label{desirables-section-quasi-coherent}

\noindent
이들은 『대수적 스택의 코호몰로지』의 절
\ref{stacks-cohomology-section-introduction}에서 정의하고 논한다.
가군의 유도 범주는 『스택의 유도 범주』의 절
\ref{stacks-perfect-section-introduction}에서 논한다.
이 장들에는 훨씬 더 많은 내용을 추가할 수 있다.



\section{평탄성과 매끄러움}
\label{desirables-section-flat-smooth}

\noindent
스킴에서 오는 평탄 전사를 갖는 것이 매끄러운 전사 조건을 대신한다는
Artin의 정리. 이 결과는 이제 『표현가능성의 판정 조건』의 정리
\ref{criteria-theorem-bootstrap}로 수록되어 있다.


\section{Artin의 표현가능성 정리}
\label{desirables-section-representability}

\noindent
이는 『Artin의 공리』의 절 \ref{artin-section-introduction}에서 논한다.
응용도 하나 있다. 『Quot 공간과 Hilbert 공간』의 정리
\ref{quot-theorem-coherent-algebraic}를 보라.
응용이 훨씬 더 많이 있어야 하며, 이 장 자체도 정리해야 한다.


\section{DM 스택은 스킴으로 유한 피복된다}
\label{desirables-section-dm-finite-cover}

\noindent
대수공간에 대한 대응하는 결과는 이미 있다. 『대수공간의 극한』의 절
\ref{spaces-limits-section-finite-cover}을 보라.
빠져 있는 것은 DM 스택과 준-DM 스택에 대한 결과이다.


\section{고유성에 관한 Martin Olsson의 논문}
\label{desirables-section-proper-parametrization}

\noindent
이 논문은 고유성의 두 개념이 같음을 증명한다. 이 논문의 첫 부분은 이제
대수적 스택에 대한 Chow의 보조정리라는 형태로 수록되어 있다.
『스택의 사상 심화』의 정리
\ref{stacks-more-morphisms-theorem-chow-finite-type}를 보라.
그 결과, 어떤 경우에는 대수적 스택의 고유성에 대한 값 판정법을 확인할 때
DVR만 사용해도 충분함을 보인다. 『스택의 사상 심화』의 절
\ref{stacks-more-morphisms-section-Noetherian-valuative-criterion}을 보라.


\section{고유 사상에 의한 연접층의 직상}
\label{desirables-section-proper-pushforward}

\noindent
이제 대수적 스택에 대한 Chow의 보조정리가 있으므로 이 주제에 관한 작업을
시작할 수 있다. 앞 절을 보라.


\section{Keel과 Mori}
\label{desirables-section-keel-mori}

\noindent
\cite{K-M}을 보라. 이들의 결과는 『스택의 사상 심화』의 절
\ref{stacks-more-morphisms-section-Keel-Mori}에 추가되었다.


\section{여기에 더 추가하기}
\label{desirables-section-add-more}

\noindent
사실 아니다. 스택 프로젝트 자체의 일부로 이 목록을 시작해서는 안 되었다!
다른 곳에 할 일 목록이 하나 있으며, 그것은 훨씬 쉽게 갱신할 수 있다.
